% !PRINTED none

\begin{center}
\gesperrt{Zweiundzwanzigster Abschnitt.}

\gesperrt{Abel'sche Körper und Kreistheilungskörper.}

\rule{3cm}{0.4pt}

\S.\,179.

\gesperrt{Einfache Abel'sche Körper.}
\end{center}

Wir wenden uns nun zu einer der interessantesten Anwendungen der Theorie der algebraischen Zahlen.

\begingroup
\renewcommand{\thefootnote}{\arabic{footnote})}
Es handelt sich um den Beweis des allgemeinen Satzes\footnote{\gesperrt{Kronecker} hat diesen Satz zuerst ausgesprochen, aber keinen vollständigen Beweis dafür veröffentlicht. Den ersten Beweis hat der Verfasser des vorliegenden Werkes in Bd. 8 der Acta Mathematica bekannt gemacht (1886). In neuester Zeit ist ein zweiter Beweis von \gesperrt{Hilbert} veröffentlicht, der sich auf die im neunzehnten Abschnitte entwickelten Sätze stützt (Nachrichten d. Gesellschaft d. Wissenschaften zu Göttingen, 1896, Heft 1).}:
\endgroup

I. \gesperrt{Alle im absoluten Rationalitätsbereich Abel'schen Zahlkörper sind Kreistheilungskörper.}

Haben wir diesen Satz bewiesen, so gewinnen die Untersuchungen des dritten Abschnittes über Kreistheilungskörper ein erhöhtes Interesse, weil damit gezeigt ist, dass auf dem dort angegebenen Wege nicht nur alle Kreistheilungskörper, sondern alle \gesperrt{Abel'schen} Körper überhaupt gefunden werden.

Es sei $\Omega = R(x)$ ein \gesperrt{Abel'scher} Körper $m^{\text{ten}}$ Grades, d. h. ein Körper, der aus allen rationalen Functionen einer Wurzel $x$ einer irreduciblen \gesperrt{Abel'schen} Gleichung $m^{\text{ten}}$ Grades besteht, wenn als Rationalitätsbereich der Körper $R$ der rationalen Zahlen betrachtet wird. Die \gesperrt{Galois'sche} Gruppe $\mathfrak{G}$ dieser Gleichung, die also eine \gesperrt{Abel'sche} Gruppe ist und denselben Grad $m$ hat, wie der Körper $\Omega$, ist auch die Gruppe des Körpers $\Omega$.

Ist $x_1$ eine nicht primitive Zahl aus $\Omega$, so ist $\Omega_1 = R(x_1)$ gleichfalls ein \gesperrt{Abel'scher} Körper, den wir einen echten Theiler von $\Omega$ nennen und dessen Grad ein echter Theiler des Grades von $\Omega$ ist. Denn ist $\mathfrak{A}$ die Gruppe, zu der die Zahl $x_1$ gehört,
%CONT>>
